\documentclass[10pt,a4paper]{amsart}
\usepackage[margin=19mm]{geometry}
\usepackage{amsmath,amssymb,mathtools}
\usepackage[scr=rsfs,cal=euler]{mathalfa}
\usepackage{xfrac}
\usepackage[unicode,hidelinks,bookmarks=false]{hyperref}
\newcommand{\ie}{i.e.\ }
\newcommand{\cf}{cf.\ }
\newcommand{\resp}{respectively\ }
\newcommand{\Subsection}{\subsection}
\providecommand{\nobreakdash}{\nobreak\hbox{-}}
\setlength{\emergencystretch}{2em}
\multlinegap0pt
\usepackage{bm}
\usepackage{bbm}
\usepackage{dsfont}
\usepackage{xspace}
\usepackage{cancel}
\usepackage{mathtools}
\usepackage{amsthm}
\makeatletter
\@ifundefined{theorem}{\newtheorem{theorem}{Theorem}}{}
\@ifundefined{lemma}{\newtheorem{lemma}[theorem]{Lemma}}{}
\makeatother
\DeclareMathOperator{\Gal}{Gal}
\DeclareMathOperator{\PSL}{PSL}
\DeclareMathOperator{\ed}{ed}
\newcommand{\Pb}{\mathbb{P}}
\newcommand{\into}{\hookrightarrow}
\newcommand{\tC}{\widetilde{C}}
\title[Essential dimension relative to branched covers of degree at]{Essential dimension relative to branched covers of degree at most n}
\author{Benson Farb, Jesse Wolfson}
\date{Source: arXiv:2510.22786v1 (CC BY 4.0)}
\begin{document}
\maketitle

\noindent\emph{Source and licence.} Essential dimension relative to branched covers of degree at most n. Version arXiv:2510.22786v1 (CC BY 4.0), distributed under the Creative Commons Attribution 4.0 International licence, as recorded in the arXiv licence metadata for this exact version and shown on the abstract page https://arxiv.org/abs/2510.22786v1. Full rights notice: licenses/arxiv-2510.22786-CC-BY-4.0.txt.\par\smallskip

\noindent\textit{Editorial scope.} Counted unit: the Lemma whose source label is \texttt{l:irred} in the article (source0.tex lines 425--433, source label \texttt{l:irred}), delivered together with the article's own complete original proof of it (source0.tex lines 434--460, one \texttt{proof} environment of 27 lines). No mathematics is written, repaired or re-derived: statement and proof are the article's own text line for line. Required context retained uncounted: l:lur (lines 407--409). Every definition, display and label of those lines is retained unchanged. Omitted from this package and not redistributed: 1--406, 410--424, 461--760. Nothing inside a retained excerpt is omitted or altered: every retained body is delivered exactly as the article's lines are written, so no omission receipt of any kind accompanies this package. Citations: the retained excerpts carry no citation command at all, so no bibliography entry of the article is transcribed and no reference is redistributed. Reference adaptation: none was needed and none was made; every reference inside the delivered unit resolves inside it, so no unresolved reference or citation is produced. Printed slips kept literal: no printed word, symbol, hypothesis or number of the counted statement or of its proof was altered. Layout and class: the article's own class is \texttt{amsart} with packages including amssymb, amsmath, amsthm, amsfonts, verbatim, tikz, amsrefs, xy, hyperref, geometry, todonotes, qtree; the wrapper is the lane's pinned \texttt{amsart} editorial wrapper at 10pt on A4, which restates the theorem-like environments the retained text needs and adds the article's own macro definitions that the retained text needs (\texttt{\textbackslash Gal}, \texttt{\textbackslash PSL}, \texttt{\textbackslash Pb}, \texttt{\textbackslash ed}, \texttt{\textbackslash into}, \texttt{\textbackslash tC}), with extra wrapper packages (bm, bbm, dsfont, xspace, cancel, mathtools). The delivered numbering is the unit's own. Assets: no retained span contains a figure environment, a graphics-inclusion command, a PDF-page inclusion, a table or a TikZ picture; no image file is copied into this package, the package declares no asset dependency and no image file is redistributed with it. Licence: version arXiv:2510.22786v1 of this article is distributed under the Creative Commons Attribution 4.0 International licence, as recorded in the arXiv licence metadata for this exact version and shown on the abstract page https://arxiv.org/abs/2510.22786v1. A full rights notice is delivered at licenses/arxiv-2510.22786-CC-BY-4.0.txt. Characters: every retained excerpt is pure ASCII, so the delivered engine, XeLaTeX with its default Latin Modern text font, reports no missing character.
\begin{lemma}\label{l:lur}
	Let $C$ be an irreducible curve. Suppose there exists a dominant map $H\to C$ and a degree $n$ rational function $h\colon H\to \Pb^1$. Then $C$ has a degree $n$ rational function $f\colon C\to \Pb^1$.
\end{lemma}

\begin{lemma}\label{l:irred}
	Let $G$ be a finite group. Let $k$ be a perfect field. Suppose that:
    \begin{enumerate}
        \item $\ed_k(G;\le n)=1$,
        \item $G$ has no proper subgroup of index at most $n$, and
        \item $G$ contains a subgroup $M\subset G$ such that $M\into \PSL_2(k)$ and $|M|>n$.
    \end{enumerate}
    Then there exists a smooth, irreducible, projective faithful $G$-curve $\tilde{C}$ with a degree $m$ rational function $f\colon\tC\to\Pb^1$ for some $m\le n$.  
\end{lemma}

\begin{proof}  We work throughout in the birational category, i.e. the category of varieties and rational maps.  

Let $V$ be a faithful representation of $G$, viewed as a linear variety. 
By assumption, there exists a branched cover $\pi\colon E\dashrightarrow V/G,$ of degree $\leq n$ 
such that $\pi^*(V \rightarrow V/G)$ arises (rationally) by pullback from a $G$-cover of smooth projective curves  
$\tilde C \rightarrow C.$
Extend the inclusion of function fields $\kappa(V/G) \rightarrow \kappa(E)$ to an inclusion of separably closed fields 
$\bar \kappa(V/G) \rightarrow \bar \kappa(E),$ and consider the maps 
$$ \Gal(\bar \kappa(E)/ \kappa(E)) \rightarrow \Gal(\bar\kappa(V/G) \rightarrow \kappa(V/G)) \rightarrow G.$$ 
corresponding to $V \rightarrow V/G.$ The second map is surjective, and the image of the composite map 
has index $\leq n,$ as $\pi$ has degree $\leq n.$ By our assumption on $G$, we conclude that the composite map is surjective, and hence $\pi^*(V) \rightarrow E$ 
is an irreducible $G$-cover. In particular, $\tilde C$ is irreducible.

Next we remark that since $V \rightarrow V/G$ is $G$-versal, it is $M$-versal. 
Indeed, if $X$ is any $M$-variety, one may apply $G$-versality to $(X \times G)/M.$ 
Thus there is a rational, $M$-equivariant map $\Pb^1 \rightarrow V,$ corresponding to 
$Y: = \Pb^1/M \rightarrow V/G.$ Let $Y_E = \pi^*(Y),$ 
let $\kappa(Y_E) \supset \kappa(Y),$ be the function field at some generic point of $Y_E,$ 
and let $\bar \kappa(Y_E)$ be a separable closure of $\kappa(Y_E).$ Consider the composite 
$$ \Gal(\bar \kappa(Y_E)/ \kappa(Y_E)) \rightarrow \Gal(\bar \kappa(Y_E)/ \kappa(Y)) \rightarrow M.$$ 
The second map is surjective, and the image of the composite has index $\leq n.$ In particular, this image is 
non-trivial, as $|M| > n.$ Hence the composite map 
$$ Y_E \rightarrow E \rightarrow C $$
is non-constant, and so the corresponding $M$-equivariant map $\tilde Y_E : \pi^*(\Pb^1) \rightarrow \tilde C$ 
is nonconstant. As $\tilde Y_E \rightarrow \Pb^1$ has degree $\leq n,$ 
we conclude by Lemma~\ref{l:lur}, that $\tC$ admits a map $h\colon \tC\to\Pb^1$ of degree $\le n$.
	\end{proof}

\end{document}
